\documentclass[11pt]{article}
\usepackage{amsmath,amssymb,amsthm}
\newtheorem{theorem}{Theorem}
\title{A confluent density test for local histogram tangent sheets}
\author{}\date{}
\begin{document}\maketitle
\paragraph{Status and scope.}
This is a complete written compactness argument, not a Lean theorem.
It concerns the first-order tangent sheets of
\texttt{local-histogram-zero-branches.tex} at a fixed finite collection
of pole profiles. It does not bound the original head model, and it
does not give a quantitative threshold at a polynomially small pole
scale. In particular it does not exclude the balanced construction
at a scale chosen separately for each number of heads.

Fix real numbers $x_{lh}$, with $1\leq l\leq g$ and
$1\leq h\leq H$. For $v=(a,c)\in\mathbb R^{2H}$ define
\[
 M_{l0}(v)=\sum_h a_h,\qquad
 M_{lj}(v)=\sum_h\bigl(a_hx_{lh}^{j}
                  +j c_hx_{lh}^{j-1}\bigr)\quad(j\geq1).
                                                               \tag{1}
\]
Let $I=[u,v_0]\subset(0,\infty)$ be a fixed compact interval.
For small positive $\epsilon$ put
\[
 T_l(\epsilon;a,c;t)=\sum_h\left(
       \frac{a_h}{t-\epsilon x_{lh}}
       +\frac{\epsilon c_h}{(t-\epsilon x_{lh})^2}\right).
                                                               \tag{2}
\]
The factor $\epsilon$ in the second term does not restrict its
coefficient space for any fixed positive $\epsilon$.

\begin{theorem}[An injective moment cutoff]
Suppose the linear map
\[
 v\longmapsto\bigl(M_{lj}(v):1\leq l\leq g,\ 0\leq j\leq L\bigr)
                                                               \tag{3}
\]
is injective. There is $\epsilon_0>0$ such that, whenever
$0<\epsilon<\epsilon_0$ and $(a,c)\ne0$, at least one function
$T_l(\epsilon;a,c;\cdot)$ has at most $L$ distinct zeros in $I$.
Consequently, on any increasing list of queries in $I$, its strict
sign word has at most $L$ sign changes.
\end{theorem}
\begin{proof}
Injectivity in finite dimension gives a constant $C$ with
\[
 \|v\|\leq C\max_{l,j\leq L}|M_{lj}(v)|.                    \tag{4}
\]
For nonzero $v$ normalize by
$\alpha=\max_{l,j\leq L}\epsilon^j|M_{lj}(v)|>0$.
For $0<\epsilon\leq1$, (4) implies
$\|v/\alpha\|\leq C\epsilon^{-L}$.
The geometric series and its derivative give
\[
 T_l(\epsilon;v/\alpha;t)
   =\sum_{j=0}^L\epsilon^jM_{lj}(v/\alpha)t^{-j-1}
       +E_l(\epsilon;t),                                    \tag{5}
\]
where $\|E_l(\epsilon;\cdot)\|_{C^L(I)}=O(\epsilon)$,
uniformly in the normalized vector. Indeed the unnormalized
remainder, including each of its first $L$ derivatives in $t$, is
$O(\epsilon^{L+1}\|v/\alpha\|)$ on $I$. All pole profiles are
fixed and $u>0$.

If the assertion were false, there would be $\epsilon_n\to0$ and
normalized vectors for which every group has at least $L+1$
distinct zeros in $I$. The finitely many coefficients
$\epsilon_n^jM_{lj}(v_n/\alpha_n)$ lie in $[-1,1]$ and their
maximum absolute value is one. Pass to a convergent subsequence.
For at least one group $l$, multiplication of (5) by $t^{L+1}$
then converges in $C^L(I)$ to a nonzero polynomial $P$ of degree
$d\leq L$. Its $d$th derivative is a nonzero constant, so for
large $n$ the $d$th derivative of
$t^{L+1}T_l(\epsilon_n;v_n/\alpha_n;t)$ has no zero in $I$.
Rolle's theorem implies that this function has at most $d$ distinct
zeros in $I$, a contradiction. Multiplication by $t^{L+1}>0$
preserves zeros and signs. Finally every adjacent strict sign change
gives a distinct zero by the intermediate value theorem.
\end{proof}

\paragraph{Finite families and sign density.}
The same conclusion holds uniformly over any finite family of profiles
for which (3) is injective with the same cutoff $L$: take the minimum
of the finitely many positive scale bounds. If there are $q$ queries
in each group, the proportion of strict sign words attained by the
union of all these tangent sheets is therefore at most
\[
 g\,2^{1-q}\sum_{j=0}^{\min(L,q-1)}\binom{q-1}{j}.            \tag{6}
\]
This is a union bound over the group having few sign changes. It is
independent of the number of sheets. A length-$q$ word with exactly
$j$ sign changes is specified by its initial sign and its $j$
change locations, which proves the count used in (6).
For example, if $L\leq(q-1)/4$, the usual binomial entropy estimate
makes (6) at most $g\,2^{-c(q-1)}$ for an absolute $c>0$.

For the histogram sheets take
$x_{lh}=\pi(h)+\rho_l(h)$. The derivative in the previous note is
exactly (2), after writing its pole perturbation as $\epsilon c$.
The theorem is applicable only after the stated moment rank has been
checked or proved for the chosen branches. Branch counting does not
establish that hypothesis, and no generic rank assertion is needed
for the theorem.

\paragraph{Why this does not settle a balanced finite-scale decoder.}
Neither (4) nor the compactness proof gives a scale bound uniform in
$H$, the profiles, or the query interval. Inverting a moment matrix
can introduce very large constants. The argument therefore does not
show that a whole-cube admissible scale of order
$1/(H^2\log H)$ lies below $\epsilon_0$. A proof at such a finite
scale, a family with delayed moment rank growth, or query locations
that themselves degenerate with $\epsilon$ remains outside this
audit. Nor does this first-order statement exclude nonlinear
perturbations of the zero branches.

\paragraph{Exact checks.}
The companion program verifies the moment coefficients of (2),
the injectivity cutoff for the three saved four-head branches,
and several higher-group examples. These are rank diagnostics,
not numerical estimates of $\epsilon_0$ or sign universality.
\end{document}
